% ===========================================================================
\chapter{Camada de Apresentação REST e Transferência de Dados (Controller \& DTO Layer)}
\label{cap_camada_apresentacao}
% ===========================================================================

A camada de apresentação REST do SIGEA-GTT é responsável por expor os endpoints da aplicação sob o protocolo HTTP, realizando o roteamento, a desserialização de requisições, a validação declarativa de payloads e a padronização de respostas JSON. Para manter a independência arquitetural preconizada pelo padrão DTO (\textit{Data Transfer Object}) \cite{fowler2002}, nenhuma entidade JPA de persistência é exposta diretamente aos consumidores da API.

% ===========================================================================
\section{Diagrama de Classes da Camada de Apresentação REST}
\label{sec_diagrama_classes_rest}
% ===========================================================================

A Figura \ref{fig_diagrama_classes_rest} detalha os controladores RESTful, o tratamento global de exceções e o envelopamento uniforme de respostas:

\begin{figure}[!ht]
\centering
\caption{Diagrama de Classes da Camada de Apresentação REST e DTOs}
\label{fig_diagrama_classes_rest}
\resizebox{0.96\textwidth}{!}{%
\begin{tikzpicture}[
    ctrl_box/.style={
        rectangle split,
        rectangle split parts=3,
        rectangle split part fill={clinical-500!20, white, white},
        draw=clinical-700,
        thick,
        rounded corners=2pt,
        font=\tiny,
        minimum width=5.2cm,
        text width=5.0cm,
        align=left
    },
    dto_box/.style={
        rectangle split,
        rectangle split parts=2,
        rectangle split part fill={teal!15, white},
        draw=teal!80!black,
        thick,
        rounded corners=2pt,
        font=\tiny,
        minimum width=4.8cm,
        text width=4.6cm,
        align=left
    },
    ex_box/.style={
        rectangle split,
        rectangle split parts=2,
        rectangle split part fill={danger-red!15, white},
        draw=danger-red!80!black,
        thick,
        rounded corners=2pt,
        font=\tiny,
        minimum width=5.2cm,
        text width=5.0cm,
        align=left
    },
    dep/.style={
        draw=clinical-900,
        thick,
        dashed,
        -Latex
    }
]

    % ====================================================
    % CONTROLADORES REST
    % ====================================================
    \node[ctrl_box] (auth_ctrl) at (-3.8, 4.5) {
        \centering\textbf{\scriptsize <<RestController>>\\AuthController}
        \nodepart{two}
        - authService : AuthService
        \nodepart{three}
        + login(LoginRequestDTO) : ResponseEntity\\
        + firstLogin(FirstLoginDTO) : ResponseEntity\\
        + changePassword(ChangePasswordDTO) : ResponseEntity
    };

    \node[ctrl_box] (user_ctrl) at (3.8, 4.5) {
        \centering\textbf{\scriptsize <<RestController>>\\UserController}
        \nodepart{two}
        - userService : UserService
        \nodepart{three}
        + getUsers(Pageable, UserRole) : ResponseEntity\\
        + getUserById(UUID) : ResponseEntity\\
        + createUser(UserCreateDTO) : ResponseEntity\\
        + updateUser(UUID, UserUpdateDTO) : ResponseEntity
    };

    \node[ctrl_box] (class_ctrl) at (-3.8, 0.0) {
        \centering\textbf{\scriptsize <<RestController>>\\AcademicClassController}
        \nodepart{two}
        - classService : AcademicClassService
        \nodepart{three}
        + createClass(ClassCreateDTO) : ResponseEntity\\
        + enrollStudents(UUID, Set<UUID>) : ResponseEntity\\
        + closeClass(UUID) : ResponseEntity\\
        + getProfessorClasses(Pageable) : ResponseEntity
    };

    \node[ctrl_box] (subm_ctrl) at (3.8, 0.0) {
        \centering\textbf{\scriptsize <<RestController>>\\ActivitySubmissionController}
        \nodepart{two}
        - submissionService : ActivitySubmissionService
        \nodepart{three}
        + submitActivity(UUID, SubmissionCreateDTO) : ResponseEntity\\
        + gradeSubmission(UUID, SubmissionGradeDTO) : ResponseEntity\\
        + getSubmissionsByActivity(UUID) : ResponseEntity
    };

    % ====================================================
    % ENVELOPADORES DE RESPOSTA E TRATAMENTO DE ERROS
    % ====================================================
    \node[dto_box] (api_resp) at (-3.8, -4.5) {
        \centering\textbf{\scriptsize <<DTO>>\\ApiResponse<T>}
        \nodepart{two}
        + success : boolean\\
        + message : String\\
        + data : T\\
        + timestamp : Instant
    };

    \node[ex_box] (ex_handler) at (3.8, -4.5) {
        \centering\textbf{\scriptsize <<RestControllerAdvice>>\\GlobalExceptionHandler}
        \nodepart{two}
        + handleBusinessException(BusinessException) : ResponseEntity\\
        + handleNotFound(ResourceNotFoundException) : ResponseEntity\\
        + handleValidation(MethodArgumentNotValidException) : ResponseEntity
    };

    % Dependências
    \draw[dep] (auth_ctrl) -- (api_resp);
    \draw[dep] (user_ctrl) -- (api_resp);
    \draw[dep] (class_ctrl) -- (api_resp);
    \draw[dep] (subm_ctrl) -- (api_resp);
    \draw[dep] (ex_handler) -- (api_resp);

\end{tikzpicture}
}
\fonte{Elaborado pelo autor (2026).}
\end{figure}

% ===========================================================================
\section{Padrão DTO e Mapeamento Estrutural com MapStruct}
\label{sec_padrao_dto_mapstruct}
% ===========================================================================

Para garantir segregação estrita de responsabilidades, os DTOs do SIGEA-GTT são divididos em objetos de comando (\textit{Request DTOs}) e objetos de projeção (\textit{Response DTOs}):

\begin{itemize}
    \item \textbf{DTOs de Entrada (\textit{Request})}: Incorporam anotações do \textit{Bean Validation} (\texttt{@NotBlank}, \texttt{@Size}, \texttt{@Email}, \texttt{@UfsEmail}), garantindo que dados inválidos sejam bloqueados na camada de transporte antes de atingir os serviços;
    \item \textbf{DTOs de Saída (\textit{Response})}: Ocultam campos sensíveis como hashes de senha, projetando apenas atributos necessários para renderização na interface;
    \item \textbf{Mappers Estruturais (MapStruct)}: As interfaces \texttt{UserMapper}, \texttt{GttModuleMapper}, \texttt{GttTriggerMapper} e \texttt{HarmSeverityMapper} são processadas em tempo de compilação, gerando código Java puro de alta performance para cópia de propriedades entre entidades e DTOs, evitando o custo computacional de reflexão em tempo de execução.
\end{itemize}

% ===========================================================================
\section{Tratamento Global de Exceções e Resposta Padronizada}
\label{sec_tratamento_excecoes}
% ===========================================================================

Toda resposta emitida pela API REST obedece ao contrato genérico envelopado pela classe \texttt{ApiResponse<T>}:
\begin{itemize}
    \item Em caso de sucesso, o atributo \texttt{success} é \texttt{true}, \texttt{message} descreve a operação concluída e \texttt{data} contém o payload do DTO tipado;
    \item Em caso de falha, o interceptador global (\texttt{GlobalExceptionHandler}, anotado com \texttt{@RestControllerAdvice}) captura a exceção da camada de serviço, mapeia o código HTTP pertinente (ex.: 400, 401, 403, 404, 409) e devolve um \texttt{ErrorResponse} estruturado sem expor detalhes internos de infraestrutura.
\end{itemize}
